% Plain-English reference for every standard-library facility used in the book.

\clearpage
\phantomsection
\section*{STL Guide for This Book}
\addcontentsline{toc}{section}{STL Guide for This Book}
\markboth{STL Guide for This Book}{C++ Standard Library}

The C++ Standard Library gives us ready-made containers and algorithms. The
name \textbf{STL} is often used for this collection, especially for containers,
iterators, and algorithms. This guide covers the standard-library tools used
by the 75 solutions in this book.

\begin{conceptbox}{Version used by the book}
Every complete solution in this edition compiles as C++11, C++17, and C++23.
The code deliberately avoids newer-only shortcuts when a clear C++11 form is
available. This means the same solution can be pasted into an older interview
environment without changing the algorithm.
\end{conceptbox}

\subsection{The common setup}

Each solution includes the headers it needs and then uses the standard
namespace. For example:

\begin{lstlisting}
#include <algorithm>
#include <string>
#include <vector>

using namespace std;

vector<int> values{4, 1, 7};
sort(values.begin(), values.end());
\end{lstlisting}

After \texttt{using namespace std;}, write \texttt{vector}, \texttt{string},
and \texttt{sort} instead of \texttt{std::vector}, \texttt{std::string}, and
\texttt{std::sort}. This is convenient in short coding-interview files. In a
large project or a header file, many teams prefer the explicit \texttt{std::}
form because it reduces the chance of two libraries using the same name.

\Needspace{25\baselineskip}
\subsection{C++11, C++17, and C++23}

\begin{revisionbox}{What changes for this book?}
\small
\setlength{\tabcolsep}{3pt}
\begin{tabularx}{\linewidth}{@{}>{\bfseries}p{1.5cm}X X@{}}
\toprule
Version & What it means here & Command-line example \\
\midrule
C++11 & The minimum version required by every solution. It provides
\texttt{vector}, hash tables, range-based loops, lambdas, \texttt{nullptr},
\texttt{move}, and the string-to-number functions used here. &
\texttt{g++ -std=c++11 file.cpp} \\
C++17 & All book code works unchanged. Features such as structured bindings,
\texttt{optional}, and \texttt{string\_view} exist, but the solutions do not
depend on them. & \texttt{g++ -std=c++17 file.cpp} \\
C++23 & All book code still works unchanged. Newer library additions may make
other programs shorter, but they do not change the algorithms taught here. &
\texttt{g++ -std=c++23 file.cpp} \\
\bottomrule
\end{tabularx}
\end{revisionbox}

The important rule is simple: choose the version accepted by the judge. If the
judge does not say, C++17 is a common practical choice. The solutions in this
book remain valid even when the judge selects C++11.

\Needspace{23\baselineskip}
\subsection{Header guide}

Only include the headers needed by the current solution.

\begin{revisionbox}{Headers used in the 75 solutions}
\footnotesize
\renewcommand{\arraystretch}{1.12}
\setlength{\tabcolsep}{3pt}
\begin{tabularx}{\linewidth}{@{}>{\raggedright\arraybackslash}p{3.0cm}
  >{\raggedright\arraybackslash}p{4.0cm}X@{}}
\toprule
Header & Main names used & Why it is used \\
\midrule
\texttt{<algorithm>} & \texttt{sort}, \texttt{min}, \texttt{max}, \texttt{reverse},
\texttt{swap} & Ordering and small reusable algorithms \\
\texttt{<array>} & \texttt{array} & Fixed-size frequency tables \\
\texttt{<cctype>} & \texttt{isalnum}, \texttt{tolower} & Character checks and conversion \\
\texttt{<cstdint>} & \texttt{uint32\_t} & An integer with an exact bit width \\
\texttt{<functional>} & \texttt{greater<int>} & Changes a priority queue into a min heap \\
\texttt{<limits>} & \texttt{numeric\_limits<int>} & Safe smallest or largest starting values \\
\texttt{<queue>} & \texttt{queue}, \texttt{priority\_queue} & BFS queues and heaps \\
\texttt{<sstream>} & \texttt{stringstream} & Reading values from a string \\
\texttt{<stack>} & \texttt{stack} & Last-in, first-out processing \\
\texttt{<string>} & \texttt{string}, \texttt{to\_string}, \texttt{stoi},
\texttt{stoull} & Text storage and number conversion \\
\texttt{<unordered\_map>} & \texttt{unordered\_map} & Key-to-value lookup \\
\texttt{<unordered\_set>} & \texttt{unordered\_set} & Fast membership tests \\
\texttt{<utility>} & \texttt{pair}, \texttt{move} & Two-value records and moving containers \\
\texttt{<vector>} & \texttt{vector} & Resizable arrays and most returned lists \\
\bottomrule
\end{tabularx}
\end{revisionbox}

\subsection{Vector and array functions}

\begin{revisionbox}{Resizable and fixed-size arrays}
\small
\setlength{\tabcolsep}{3pt}
\begin{tabularx}{\linewidth}{@{}>{\ttfamily}p{3.1cm}X>{\raggedright\arraybackslash}p{2.2cm}@{}}
\toprule
Operation & Meaning & Usual cost \\
\midrule
v[i] & Read or change the value at index \texttt{i}. & $O(1)$ \\
v.size() & Number of stored values. & $O(1)$ \\
v.empty() & True when the container has no values. & $O(1)$ \\
v.push\_back(x) & Add \texttt{x} at the end. & $O(1)$ average \\
v.back() & Read the final value. & $O(1)$ \\
v.reserve(n) & Reserve room before repeated appends. & $O(n)$ when storage moves \\
v.clear() & Remove all values. & Linear in stored objects \\
v.begin(), v.end() & Iterator range covering every value. & $O(1)$ \\
a.fill(x) & Put \texttt{x} in every cell of a fixed \texttt{array}. & $O(n)$ \\
\bottomrule
\end{tabularx}
\end{revisionbox}

\begin{lstlisting}
#include <array>
#include <vector>

using namespace std;

vector<int> values;
values.reserve(3);
values.push_back(4);
values.push_back(7);

array<int, 26> count{};
count.fill(0);
\end{lstlisting}

Use \texttt{vector} when the size changes. Use \texttt{array<T,N>} when the
size is known at compile time, such as the 26 lowercase English letters.

\Needspace{23\baselineskip}
\subsection{String functions}

\begin{revisionbox}{Text operations used in the book}
\small
\setlength{\tabcolsep}{3pt}
\begin{tabularx}{\linewidth}{@{}>{\raggedright\arraybackslash}p{3.6cm}X@{}}
\toprule
Operation & Meaning and important note \\
\midrule
\texttt{s[i]} & Character at index \texttt{i}; check the index first. \\
\texttt{s.size()} & Number of characters; the return type is unsigned. \\
\texttt{s.empty()} & True when the string has no characters; clearer than \texttt{s.size()==0}. \\
\texttt{substr(pos,len)} & Copy part of a string; cost is proportional to the copied text. \\
\texttt{s.find(ch,pos)} & Find the next character or substring; returns \texttt{string::npos} if absent. \\
\texttt{s.push\_back(ch)} & Add one character; constant time on average. \\
\texttt{to\_string(x)} & Convert a number to a string; available in C++11. \\
\texttt{stoi(s)} & Convert text to \texttt{int}; may throw for invalid text. \\
\texttt{stoull(s)} & Convert text to \texttt{unsigned long long}; useful for stored lengths. \\
\bottomrule
\end{tabularx}
\end{revisionbox}

\begin{lstlisting}
#include <string>

using namespace std;

string encoded = "4#code";
size_t mark = encoded.find('#');
int length = stoi(encoded.substr(0, mark));
string word = encoded.substr(mark + 1, length);
\end{lstlisting}

\subsection{Hash map and hash set functions}

Both containers use hashing. Average lookup, insertion, and removal are
$O(1)$, although the worst case can be $O(n)$.

\begin{revisionbox}{Fast lookup operations}
\small
\setlength{\tabcolsep}{3pt}
\begin{tabularx}{\linewidth}{@{}>{\raggedright\arraybackslash}p{3.5cm}X@{}}
\toprule
Operation & Meaning \\
\midrule
\texttt{table[key]} & Read or create the value stored for \texttt{key}. \\
\texttt{table.find(key)} & Return an iterator to the key, or \texttt{table.end()}. \\
\texttt{table.count(key)} & Return 0 or 1 for the unique-key containers used here. \\
\texttt{table.insert(x)} & Add a key or a key-value pair. \\
\texttt{table.clear()} & Remove all entries. \\
\texttt{for (auto\& entry : table)} & Visit every stored entry; the order is unspecified. \\
\bottomrule
\end{tabularx}
\end{revisionbox}

\begin{lstlisting}
#include <unordered_map>
#include <unordered_set>

using namespace std;

unordered_map<int, int> indexOf;
indexOf[7] = 3;
if (indexOf.count(7)) {
    int index = indexOf[7];
}

unordered_set<int> seen;
seen.insert(12);
bool alreadySeen = seen.count(12) != 0;
\end{lstlisting}

Do not use \texttt{operator[]} only to test whether a map key exists: it may
create a new entry. Use \texttt{count} or \texttt{find} for a pure lookup.

\Needspace{14\baselineskip}
\subsection{Queue, stack, and priority queue}

\begin{revisionbox}{The accessible end of each structure}
\small
\setlength{\tabcolsep}{3pt}
\begin{tabularx}{\linewidth}{@{}>{\raggedright\arraybackslash\bfseries}p{2.2cm}X X@{}}
\toprule
Type & Add, read, remove & Typical use \\
\midrule
Queue & \texttt{push}, \texttt{front}, \texttt{pop} & BFS and first-in,
first-out work \\
Stack & \texttt{push}, \texttt{top}, \texttt{pop} & Nested delimiters and
last-in, first-out work \\
Max heap & \texttt{push}, \texttt{top}, \texttt{pop} & Repeatedly
take the largest value \\
Min heap & \texttt{push}, \texttt{top}, \texttt{pop} & Repeatedly take the
smallest value; the full type appears below. \\
\bottomrule
\end{tabularx}
\end{revisionbox}

\begin{lstlisting}
#include <functional>
#include <queue>
#include <vector>

using namespace std;

queue<int> bfs;
bfs.push(3);
int nextNode = bfs.front();
bfs.pop();

priority_queue<int> maxHeap;
priority_queue<int, vector<int>, greater<int>> minHeap;
\end{lstlisting}

Calling \texttt{pop()} does not return the removed value. Read \texttt{front()}
or \texttt{top()} first, then call \texttt{pop()}.

\subsection{Algorithms used in the book}

Standard algorithms normally receive a half-open range:
\texttt{[begin,end)}. The first iterator is included; the end iterator points
one position after the final element.

\begin{revisionbox}{Algorithm quick reference}
\small
\setlength{\tabcolsep}{3pt}
\begin{tabularx}{\linewidth}{@{}>{\raggedright\arraybackslash}p{3.6cm}X>{\raggedright\arraybackslash}p{2.0cm}@{}}
\toprule
Operation & Meaning & Cost \\
\midrule
\texttt{sort(range)} & Sort in increasing order. & $O(n\log n)$ \\
\texttt{sort(range,rule)} & Sort using a custom rule. & $O(n\log n)$ \\
\texttt{min(a,b)}, \texttt{max(a,b)} & Return the smaller or larger value. & $O(1)$ \\
\texttt{reverse(range)} & Reverse a range in place. & $O(n)$ \\
\texttt{swap(a,b)} & Exchange two values. & Usually $O(1)$ \\
\bottomrule
\end{tabularx}
\end{revisionbox}

The following comparison works in C++11 because its parameter types are
written explicitly:

\begin{lstlisting}
#include <algorithm>
#include <vector>

using namespace std;

vector<vector<int>> intervals{{1, 3}, {2, 4}, {3, 5}};
sort(intervals.begin(), intervals.end(),
     [](const vector<int>& a, const vector<int>& b) {
         return a[1] < b[1];
     });
\end{lstlisting}

Writing \texttt{const auto\&} as a lambda parameter is shorter, but generic
lambda parameters require C++14 or newer.

\Needspace{22\baselineskip}
\subsection{Numbers, characters, streams, and moves}

\begin{revisionbox}{Other library tools used by the solutions}
\small
\setlength{\tabcolsep}{3pt}
\begin{tabularx}{\linewidth}{@{}>{\raggedright\arraybackslash\bfseries}p{3.0cm}X@{}}
\toprule
Name & Purpose \\
\midrule
Smallest integer & \texttt{numeric\_limits<int>::min()} is useful when every
valid answer may be negative. \\
Largest integer & \texttt{numeric\_limits<int>::max()} is useful as an initial
upper bound. \\
Exact 32-bit integer & \texttt{uint32\_t} is unsigned and has exactly 32 bits
when the implementation provides it. \\
Character check & \texttt{isalnum(ch)} tests whether a character is a letter
or digit. \\
Lowercase conversion & \texttt{tolower(ch)} produces the lowercase form of a
character. \\
String stream & \texttt{stringstream} reads formatted values from a string. \\
Move & \texttt{move(value)} allows a container's stored data to be transferred
instead of copied when possible. \\
Pair & \texttt{pair<A,B>} stores two related values in \texttt{first} and
\texttt{second}. \\
\bottomrule
\end{tabularx}
\end{revisionbox}

When calling \texttt{isalnum} or \texttt{tolower}, cast a possibly negative
\texttt{char} to \texttt{unsigned char} first. This avoids undefined behavior
for characters outside the basic ASCII range.

\Needspace{25\baselineskip}
\subsection{Newer shortcuts not required here}

The following features are useful, but using them would raise the minimum
language version. The right column shows the portable form used in this book.

\begin{revisionbox}{Version-sensitive alternatives}
\small
\setlength{\tabcolsep}{3pt}
\begin{tabularx}{\linewidth}{@{}>{\raggedright\arraybackslash\bfseries}p{2.7cm}
  >{\centering\arraybackslash}p{1.5cm}
  >{\raggedright\arraybackslash}X@{}}
\toprule
Newer feature & First version & Form used in this book \\
\midrule
Generic lambda & C++14 & Replace an \texttt{auto} parameter with the full type. \\
Structured bindings & C++17 & Use \texttt{pair.first}, \texttt{pair.second},
or an ordinary loop variable. \\
Views and optional values & C++17 & Use \texttt{string}, indexes, pointers, or a
separate Boolean value. \\
Membership check & C++20 & Replace \texttt{contains(key)} with \texttt{count(key)} or
\texttt{find(key)!=end()}. \\
Range sorting & C++20 & Replace \texttt{ranges::sort(values)} with
\texttt{sort(values.begin(),values.end())}. \\
Bit count & C++20 & Replace \texttt{popcount(value)} by clearing the lowest set bit in a loop. \\
Safe midpoint & C++20 & Use \texttt{left+(right-left)/2}. \\
Formatted print & C++23 & Online-judge solutions return values; they do not need
formatted printing. \\
\bottomrule
\end{tabularx}
\end{revisionbox}

\subsection{Fast checklist before submitting}

\begin{revisionbox}{STL mistakes that often cause wrong answers}
\begin{itemize}
  \item Include the correct header instead of relying on another header to
  include it by accident.
  \item Check \texttt{empty()} before calling \texttt{front()}, \texttt{back()},
  or \texttt{top()}.
  \item Remember that \texttt{pop()} removes an item but does not return it.
  \item Do not dereference \texttt{end()} after an unsuccessful \texttt{find()}.
  \item Remember that \texttt{unordered\_map} iteration order is not sorted.
  \item Use \texttt{static\_cast<int>(container.size())} when mixing a size with
  signed integer indexes.
  \item Count sorting and heap operations as $O(n\log n)$ when explaining the
  final complexity.
\end{itemize}
\end{revisionbox}

The rest of the book introduces these tools inside real problems. Return to
this guide whenever a container operation or header name is unfamiliar.